\documentclass[11pt,letterpaper]{article}
\usepackage[margin=1.2in]{geometry}
\usepackage[scaled=0.95]{helvet}
\renewcommand{\familydefault}{\sfdefault}
\usepackage[T1]{fontenc}
\usepackage{microtype}
\usepackage{xcolor}
\definecolor{softblack}{HTML}{1A1A1A}
\definecolor{mutedgray}{HTML}{555555}
\AtBeginDocument{\color{softblack}}
\setlength{\parskip}{6pt}
\setlength{\parindent}{0pt}
\pagestyle{empty}

\begin{document}
\noindent\textbf{Cover letter --- \emph{Quarterly Journal of Economics}}\\[0.4em]
\textcolor{mutedgray}{\rule{0.22\linewidth}{0.4pt}}

\bigskip
\noindent Editor\\
\textit{Quarterly Journal of Economics}

\bigskip

\noindent Dear Editor,

\bigskip

\noindent I am submitting ``The Long Shadow of Seasonality: Climate Endowments, Agricultural Pathways, and Escape from the Malthusian Trap'' for consideration at the \emph{QJE}.

The paper extends \citet{matranga2024} into the post-Neolithic period and produces three findings that I believe make it an appropriate companion to that work. The argument is that the climate dimension Matranga identified at the origin of agriculture also shapes which agricultural system a society adopts after the Neolithic, and that the selected system in turn governs the Malthusian feedback and the response to climate shocks all the way to the modern transition.

\textbf{Empirical contributions.} \emph{First}, I construct a country-level and sub-national monthly climate panel reaching back to 1421~CE by area-weighting the ModE-RA paleo-reanalysis onto HYDE 3.5 sub-units. The country-month correlation against country-level ERA5 (which I rebuild from the raw hourly fields onto HYDE boundaries) is $0.977$ in median; the annual anomaly correlation is $0.832$ over the full 59-year overlap. This is, to my knowledge, the first dataset of its kind for econometric work and substantially expands what can be tested in the climate-and-pre-industrial-growth literature.

\emph{Second}, I show that raw temperature seasonality alone, the natural Matranga measure, does \emph{not} significantly distinguish agricultural pathways at the country level once climate is measured properly ($F=0.50$, $p=0.68$ on $N=153$ countries). What does is the productive-months index---the joint structure of temperature and precipitation that determines whether rain-fed agriculture is feasible ($F=6.31$, $p<0.001$). This is a sharpening of the Matranga story, not a refutation: storage demand is generated by the joint constraint of growing-season conditions and the off-season's harshness, not by temperature variability in isolation. I view the QJE-natural framing as ``Matranga's storage-demand mechanism, sharpened and extended into the post-Neolithic.''

\emph{Third}, exploiting HYDE's sub-national grid (3{,}187 sub-units across 196 countries) and adding country fixed effects, I identify the storage-demand effect on agricultural specialisation off purely within-country variation. The within-country point estimate is 3.5$\times$ the between-country estimate, consistent with cross-country regressions being attenuated by mean-reverting institutional controls. Within-country, an extra productive month per year is associated with 6.7~pp higher crop share at 1750 ($p=0.019$ with country-clustered SE, $p=0.0001$ with two-way clustering by country and 10$^\circ$ latitude band) and 0.66 log-units higher population density ($p<10^{-3}$). This is the paper's tightest causal evidence and the move that puts the long-run climate-development literature on a within-country footing.

A 1421--1950 panel with real annual climate also surfaces two by-product findings: a within-paper signature of the demographic transition (the Malthusian density coefficient progressively vanishing across subperiods) and a sign-flip on inter-annual temperature volatility between 1750--1900 and 1900--1950 ($-0.0125$, $p < 10^{-4}$), consistent with modern monoculture being more volatility-sensitive than diversified pre-industrial systems.

The paper is single-authored and the empirical infrastructure (ModE-RA aggregation, country-level ERA5 rebuild, sub-national panels) is built from scratch. A replication package will be deposited on OSF upon acceptance. The work has not been previously published and is not under consideration elsewhere; I declare no conflicts.

\bigskip

\noindent Sincerely,\\[0.4em]
Jorge Alonso Ortiz\\[0.2em]
\small\textcolor{mutedgray}{[Affiliation]\quad\textbar\quad jorge.alonso@icloud.com}

\bigskip

\textcolor{mutedgray}{\rule{0.22\linewidth}{0.4pt}}\\[0.2em]
{\footnotesize\textcolor{mutedgray}{Suggested referees (independent of the author):}
A. Matranga (NES); O. Galor or D. Weil (Brown); Q. Ashraf (Williams); N. Voigtländer (UCLA); S. Hsiang (Stanford). I will note that the empirical strategy speaks most directly to readers familiar with Matranga (2024) and Ashraf--Galor (2011).}

\end{document}
